% Chapter 16 — Minimum Requirements. Author: NAZAROV. Budget 4 pp.
% Discharges Constitution v1.2 §14.1 and §14.2. Restates as minima the derivations of
% Ch.4 (machines, §14.1), Ch.7 (valuation P1–P6, §14.2), Ch.8 (six auditor checks, §14.2).
% No thread episodes: requirements chapters cite episodes, they do not re-run them.
\section{Minimum Requirements}\label{ch:requirements}

This chapter discharges constitution clauses C-14.1--C-14.15 (the minimum requirements for the untrusted machines and the pricing layers).

It adds nothing. Each requirement below is a consequence already derived and proved in
the body --- restated once, imperatively, as the minimum any implementation must meet, and
traceable by cross-reference to the chapter that established it. The requirements fall in
three groups: the Event Monitor and the Events Executor (\S 14.1, from
Chapter~\ref{ch:machines}); the valuation semantics (\S 14.2, from
Chapter~\ref{ch:valuation}); and the data boundary (\S 14.2, from
Chapter~\ref{ch:marketdata}). A candidate implementation is conformant when every
requirement holds; \S 16.4 states how an auditor confirms it.

\begin{principle}[What the minima secure]
The ledger's integrity holds by construction and rests on the Transaction Executor alone
(Chapter~\ref{ch:machines}). The requirements of this chapter secure the properties
integrity does \emph{not} supply: the liveness and replayability of the untrusted
machinery, and the reproducibility and boundary-closure of the pricing layers. Integrity
owes these components nothing; timeliness, reproducibility, and the closed-system property
are conditional on them.
\end{principle}

\subsection{The Event Monitor and the Events Executor}

The Event Monitor and the Events Executor sit outside the trust boundary
(Chapter~\ref{ch:machines}). No failure either can commit corrupts ledger state; every one
must degrade to a visible, overdue item rather than to silence. What they owe is liveness
and replayability, and the following seven requirements are its minimum.

\begin{enumerate}[label=\textbf{M\arabic*.}, ref=M\arabic*, leftmargin=*, itemsep=3pt]
\item \textbf{Durable watches and timers.} A watch registered for a future date or
  condition must fire when the condition is met, across restarts and failures; timer state
  must be persisted, never held only in volatile memory (Chapter~\ref{ch:machines}). A
  watch that could be lost on restart is a silent gap, which \S 14.1 forbids.
\item \textbf{At-least-once emission.} No emitted event may be lost; the Monitor must
  re-emit until the event is on the record. A duplicate emission must be harmless, made so
  by the cause-derived idempotence check at the door (Chapter~\ref{ch:machines}). Losing an
  event is unrecoverable; emitting one twice costs nothing.
\item \textbf{Acknowledged watches.} Every declared watch must be acknowledged by the
  Monitor, and the acknowledgement recorded in the unit's state, so a watch is visibly
  declared, then armed, then fired or expired. A watch unacknowledged beyond its arming
  window must stand on the record as an overdue item, never as absence
  (Chapter~\ref{ch:machines}).
\item \textbf{Triggers carry timing, not data.} A trigger must say only \emph{when};
  everything the event means must be read from the record by the contract it fires
  (Chapter~\ref{ch:machines}). A trigger carrying data would place an unrecorded input
  before the door and break reproducibility.
\item \textbf{Deterministic, replayable orchestration.} Each machine's own state must be
  reconstructible by replaying its append-only history, so a late firing produces the
  identical transaction and any party can recompute which events should have been emitted,
  and when (Chapter~\ref{ch:machines}).
\item \textbf{No write privilege.} Neither the Event Monitor nor the Events Executor may
  hold write access to the ledger. The Transaction Executor is the single writer
  (Chapter~\ref{ch:machines}); an emission or a firing is a proposal, never a commit.
\item \textbf{Obligation liveness.} Every contractual obligation must be registered with a
  deadline, a test for its discharge, and a fallback action; by its deadline it must be
  discharged, compensated, or declared defaulted. A duty that fails to fire must degrade to
  a visible, overdue item, never to silence. This is the liveness counterpart of the value
  sufficiency obligation of Chapter~\ref{ch:collateral} and Chapter~\ref{ch:invariants}.
\end{enumerate}

\subsection{Valuation semantics}

Net Asset Value is a projection of the log, folded and priced; valuation adds no store
(Chapter~\ref{ch:valuation}). Chapter~7 derived six requirements (P1--P6) on the pricing
stack and the pricing data layer as consequences of that fact. They are the minima any
pricing layer must meet.

\begin{enumerate}[label=\textbf{V\arabic*.}, ref=V\arabic*, leftmargin=*, itemsep=3pt]
\item \textbf{The ledger records pricing outputs and never runs a model.} A pricing model
  must run outside the ledger; its output must re-enter only as a recorded observation
  carrying sufficient provenance (Chapter~\ref{ch:valuation}, P1). Reproducing a valuation
  from recorded prices is then unconditional; reproducing the model's number requires the
  model itself, whose retention is a governance matter outside scope
  (Chapter~\ref{ch:scope}).
\item \textbf{Valuation is state-parameterised, and the state comes from the ledger.} The
  price function must read the unit's lifecycle state from UnitStatus and PositionState
  (Chapter~\ref{ch:homes}), never from a store the pricing layer keeps of its own
  (Chapter~\ref{ch:valuation}, P2). Value that turns on state turns on state the ledger
  alone authors.
\item \textbf{Valuations are reproducible from recorded inputs.} Every NAV must be
  recomputable to the minor unit from the composition --- a fold of the log --- and prices
  identified on the record (Chapter~\ref{ch:valuation}, P3). No valuation may depend on an
  input that is not on the record.
\item \textbf{Prices and quantities are never mixed across an event.} A price and a
  quantity may be combined only when both refer to the same state of the instrument ---
  consistency of reference (Chapter~\ref{ch:valuation}, P4; catalogued in
  Chapter~\ref{ch:invariants}). The mechanism that enforces this at ingestion is
  requirement~B5.
\item \textbf{Two price vectors are both projections over one composition.} Mark-to-market
  and mark-to-mid must be two price vectors applied to the single composition the fold
  produces; positions must never diverge --- only the two money figures may
  (Chapter~\ref{ch:valuation}, P5). One position record forbids the settlement-versus-risk
  position drift that is a classical reconciliation break.
\item \textbf{Estimates are kept apart from entitlements.} An estimated or model-derived
  value must never be recorded as, or substituted for, an owed quantity: a mark is a
  projection, an entitlement is a recorded fact (Chapter~\ref{ch:valuation}, P6). The
  complementary boundary duty --- originals preserved, adjusted and derived values
  recomputed at read --- is requirement~B4.
\end{enumerate}

\subsection{The data boundary}

The boundary is the one place the closed system is not closed
(Chapter~\ref{ch:marketdata}). Chapter~8 held it watertight with two devices --- the
recorded observation and the market data operator --- and listed six checks a candidate
implementation must pass. A seventh extends the same discipline to the declared terms of a
governing agreement, whose provenance is the executed agreement rather than a market feed
(Chapter~\ref{ch:collateral}). They are the boundary minima.

\begin{enumerate}[label=\textbf{B\arabic*.}, ref=B\arabic*, leftmargin=*, itemsep=3pt]
\item \textbf{All data carries provenance.} All data on the record must carry its
  provenance --- source, observation time, and basis frame --- and no valuation may read
  data that does not (Chapter~\ref{ch:marketdata}, \S\ref{sec:obs-door}). A bare read
  against a live feed leaves no record and is unreproducible; it must be refused.
\item \textbf{Kinds are registered before they cross.} Every data kind must be
  registered, with the dimension of each component --- price, proportional, or quantity ---
  before any data of that kind crosses; data of an unregistered kind must be refused at
  the door (Chapter~\ref{ch:marketdata}, \S\ref{sec:registry}). Registration is immutable:
  a correction is a new kind plus a retirement of the old, never an edit.
\item \textbf{An operator is declared for every pairing.} For every (registered data kind,
  event kind) pair the boundary admits, a market data operator must be declared --- the
  identity operator where a kind is genuinely unaffected, but declared, not assumed. A
  pairing with no declared operator must be refused, never defaulted to identity
  (Principle~\ref{prin:sched-total}, Chapter~\ref{ch:marketdata}). This is
  adjustment-schedule totality, and it fails closed.
\item \textbf{Originals preserved, adjusted computed at read.} All observed data must be
  preserved on the record exactly as observed; every adjusted value must be recomputed at
  each read from the original under the declared operators, and every derived quantity
  recomputed from operator-adjusted inputs
  (Chapter~\ref{ch:marketdata}, \S\ref{sec:menu}--\S\ref{sec:recompose}). An original must
  never be overwritten by its adjustment.
\item \textbf{The invariance witness holds.} Across every corporate action the operator on
  the value frame and the move on the quantity must be bound so that total value is
  invariant across the event; a stale-frame read --- a new count against an old price ---
  or a wrong-side-of-ex read must be refused
  (Chapter~\ref{ch:marketdata}, \S\ref{sec:witness}). This is the ingestion-side
  enforcement of requirement~V4.
\item \textbf{Each failure mode has its determined disposition.} Each boundary failure
  --- missing data (W1), late data (W2), contradictory sources (W3), unmappable event
  (W4) --- must produce its determined, replayable disposition, a function of log state
  alone, so replay reproduces it bit-identically; none may improvise or pick silently
  (Chapter~\ref{ch:marketdata}, \S\ref{sec:failures}). In particular a single source for
  data that materially affects valuation is insufficient: such data must be drawn from
  more than one attested source, aggregated by a declared rule, with divergence beyond a
  declared threshold producing a flagged output rather than a silent pick (W3).
\item \textbf{Declared agreement terms are reconciled at the boundary.} Every declared
  term of a governing agreement --- eligibility, valuation percentage, thresholds, the
  trapping predicate, and the legal-regime bit --- is boundary data whose provenance is the
  executed agreement; it must be reconciled against the counterparty's own record, and a
  transcription error repaired by terms amendment plus an authorised compensating rebooking, never a
  silent edit (Chapter~\ref{ch:collateral}; Chapter~\ref{ch:marketdata}). The regime bit is
  the worked instance; the duty is general over every declared term.
\end{enumerate}

The boundary rests on three named trust assumptions and no other: TA-KIND, that a kind
declaration faithfully renders its publisher's quotation convention, owned by data
governance and detected by the invariance witness, source divergence, and recomputation
(Chapter~\ref{ch:marketdata}); TA-TERMS, that a declared agreement term faithfully
transcribes the executed agreement (B7), owned by the same governance and detected by
boundary reconciliation against the counterparty's own record; and TA-EXDATE, that the
recorded ex-date coheres with the record date and the settlement cycle so that a trade's
registration state at the record date agrees with its cum/ex entitlement. \emph{Owner:}
the party that records the corporate-action calendar (data governance). \emph{Violation:}
an ex trade settles before the record date --- lawfully, under a due-bills special dividend
whose ex-date falls after the payment date --- so the registered holder, the buyer, is not
the entitled party, the seller. \emph{Detection:} the entitlement routing itself, which in
that fourth quadrant raises the mirror market-claim leg from buyer to seller
(Chapter~\ref{ch:settlement}) rather than misrouting silently, together with boundary
reconciliation of the recorded ex-date against the external corporate-action calendar.
\emph{Residual:} the coherence is an external convention nowhere on the record, so a
recorded ex-date that misstates it misroutes faithfully to the record --- a perimeter
reconciliation matter, the sibling of TA-KIND and TA-TERMS. The fourth quadrant is handled
on the record either way: the mirror market-claim leg routes the distribution to the
entitled seller whether or not the convention held, so TA-EXDATE names a residual the
ledger reconciles against, not a gap in the routing. Everything else that crosses is
recorded, framed, and replayable.

\subsection{Verification}

An auditor confirms conformance by exercising each group against generated products,
events, and histories, on the executable-check surface of Chapter~\ref{ch:testability}.

The machine minima (M1--M7) are liveness and replay properties: generate restart,
duplicate-emission, and late-firing histories, and confirm that no event is lost, that
every watch is acknowledged or stands overdue, that every late firing reproduces the
identical transaction, and that neither machine writes the ledger. The valuation minima
(V1--V6) are checkable over generated positions and prices: recompute every NAV from the
record and compare to the minor unit, confirm the price function reads state only from the
ledger, and confirm that mark-to-market and mark-to-mid fold one composition. The boundary
minima (B1--B7) are checkable over generated data, corporate actions, and agreements:
confirm that no
data lacks provenance, that no unregistered kind and no undeclared pairing crosses, that
every original survives its adjustment, that the invariance witness holds across every
generated event, that each of W1--W4 replays bit-identically, and that a perturbed
agreement term is caught by boundary reconciliation and repaired by amendment plus
rebooking rather than silent edit
(Chapters~\ref{ch:testability} and~\ref{ch:invariants}).

Because every requirement here restates a derivation already proved in the body,
conformance is not an added burden but the body's own claims made executable. When they
hold, the untrusted machinery is live and replayable, the pricing layer adds no store, and
the boundary is closed --- so the conservation and reproducibility the rest of the
specification proves rest on stated, testable minima.
